%
% 3-fortsetzung.tex -- Analytische Fortsetzung
%
% (c) 2025 Prof Dr Andreas Müller
%
\section{Analytische Fortsetzung
\label{buch:analytisch:section:fortsetzung}}
\kopfrechts{Analytische Fortsetzung}
Auf den ersten Blick scheint die Beschreibung einer holomorphen
Funktion durch eine Potenzreihe einen Verlust zu beinhalten.
Eine Potenzreihe ist ja nur innerhalb des Konvergenzkreises
konvergent, während eine holomorphe Funktion einen viel grösseren
Definitionsbereich haben kann.
Das Definitionsgebiet lässt sich aber immer mit Kreisen überdecken,
innerhalb derer sich die Funktion durch eine Potenzreihe schreiben
lässt.
In den Überschneidungsgebieten dieser Kreise stimmen die Funktionen
überein.
Der in diesem Abschnitt dargestellte Prozess der analytischen
Fortsetzung ermöglicht, eine holomorphe Funktion entlang einer
Kette von Kreisgebieten zu erweitern.
So lässt sich das grösstmögliche Definitionsgebiet
für eine holomorphe Funktion finden.
Es zeigt aber zum Beispiel bei der Logarithmus-Funktion auch, dass
die Menge der komplexen Zahlen zu klein sein kann und durch eine
Überlagerung ersetzt werden muss, um eine wohldefinierte komplex
differenzierbare Funktion zu konstruieren.

\input{chapters/040-analytisch/31-definition.tex}
\input{chapters/040-analytisch/32-fortsetzung.tex}
\input{chapters/040-analytisch/33-ueberlagerungen.tex}
